% GENERATED FILE. Edit canonical lessons, then run npm run generate:books.
% canonical-source-hash: fnv1a64:d14155c238c83090
% canonical-lessons: BN-C508-boithokkhana, BN-C508-chilekotha, BN-C508-rannaghor, BN-C508-choukath, BN-C508-tangano

\chapter{Rooms of a House}
\label{ch:bn-rooms-of-a-house}
\hlchaptermodality{508}

\begin{chapteropening}
\textbf{By the end of this chapter:} \emph{I can name the rooms of a house and hang things on the walls.}

Complete the last lesson of chapter 508: I can name the rooms of a house and hang things on the walls.
\end{chapteropening}

\section[\texorpdfstring{\bn{বৈঠকখানা} (boiṭhôkkhānā)}{বৈঠকখানা (boiṭhôkkhānā)}]{\bn{বৈঠকখানা} (boiṭhôkkhānā) — a sitting room; a drawing room}

\label{lesson:BN-C508-boithokkhana}



\begin{quote}
Before the new one: say the Bengali for to plaster, then the Bengali for to hammer in.
\end{quote}

\subsection*{You'll want to know: \bn{বৈঠকখানা}}
\textbf{\bn{বৈঠকখানা}} — \emph{boiṭhôkkhānā} — \textquotedblleft{}a sitting room; a drawing room\textquotedblright{}.

\begin{quote}
\bn{দাদু} \bn{বৈঠকখানায়} \bn{বসে} \bn{পত্রিকা} \bn{পড়েন।}
\end{quote}

\begin{quote}
\emph{Dādu boiṭhôkkhānāy bôse pôtrikā pôṛen.}
\end{quote}

\begin{quote}
My grandfather sits in the drawing room and reads the magazine.
\end{quote}

\begin{cousinweb}
\textbf{\bn{বৈঠকখানা}} (\emph{boiṭhôkkhānā}) joins \bn{বৈঠক}, \textquotedblleft{}a sitting, a meeting\textquotedblright{}, usually traced to Hindi \emph{baiṭhak}, to \emph{khāna}, \textquotedblleft{}a house, a room\textquotedblright{}, which is usually traced to Persian.
\end{cousinweb}

\begin{grammarlens}[title={What you've built}]
One more word for describing what you lived and felt.
\end{grammarlens}

\subsection*{Your turn}
\begin{itemize}
\raggedright
  \item \emph{Say it:} \emph{boiṭhôkkhānā}
  \item \emph{Say it:} \emph{boiṭhôkkhānā}, once more
  \item \emph{Say it:} \emph{Dādu boiṭhôkkhānāy bôse pôtrikā pôṛen.} — My grandfather sits in the drawing room and reads the magazine.
  \item \emph{Recall:} say the Bengali for to plaster, then the Bengali for to hammer in, then say \emph{boiṭhôkkhānā} again
\end{itemize}

\begin{tcolorbox}[breakable,colback=teal!4,colframe=teal!35!black,title={Before you move on}]
What does \bn{বৈঠকখানা} mean? (\textquotedblleft{}A sitting room; a drawing room\textquotedblright{}.) Say it once more.
\end{tcolorbox}
\section[\texorpdfstring{\bn{চিলেকোঠা} (chilekoṭhā)}{চিলেকোঠা (chilekoṭhā)}]{\bn{চিলেকোঠা} (chilekoṭhā) — a small room on the roof; an attic}

\label{lesson:BN-C508-chilekotha}



\begin{quote}
Before the new one: say the Bengali for to hammer in, then the Bengali for a drawing room.
\end{quote}

\subsection*{You'll want to know: \bn{চিলেকোঠা}}
\textbf{\bn{চিলেকোঠা}} — \emph{chilekoṭhā} — \textquotedblleft{}a small room on the roof; an attic\textquotedblright{}.

\begin{quote}
\bn{পুরনো} \bn{খেলনাগুলো} \bn{চিলেকোঠায়} \bn{আছে।}
\end{quote}

\begin{quote}
\emph{Purôno khelnāgulo chilekoṭhāy āchhe.}
\end{quote}

\begin{quote}
The old toys are in the room on the roof.
\end{quote}

\begin{cousinweb}
\textbf{\bn{চিলেকোঠা}} (\emph{chilekoṭhā}) is the small room at the top of the stairs, on the roof. The second part is \bn{কোঠা}, \textquotedblleft{}a room\textquotedblright{}, which CDIAL lists as Bengali \emph{koṭhā}, \textquotedblleft{}a brick-built house\textquotedblright{}, under Sanskrit \emph{kōṣṭha}, \textquotedblleft{}a storeroom, an inner apartment\textquotedblright{} (CDIAL 3546); where the first part comes from is not settled.
\end{cousinweb}

\begin{grammarlens}[title={What you've built}]
One more word for describing what you lived and felt.
\end{grammarlens}

\subsection*{Your turn}
\begin{itemize}
\raggedright
  \item \emph{Say it:} \emph{chilekoṭhā}
  \item \emph{Say it:} \emph{chilekoṭhā}, once more
  \item \emph{Say it:} \emph{Purôno khelnāgulo chilekoṭhāy āchhe.} — The old toys are in the room on the roof.
  \item \emph{Recall:} say the Bengali for to hammer in, then the Bengali for a drawing room, then say \emph{chilekoṭhā} again
\end{itemize}

\begin{tcolorbox}[breakable,colback=teal!4,colframe=teal!35!black,title={Before you move on}]
What does \bn{চিলেকোঠা} mean? (\textquotedblleft{}A small room on the roof; an attic\textquotedblright{}.) Say it once more.
\end{tcolorbox}
\section[\texorpdfstring{\bn{রান্নাঘর} (rānnāghôr)}{রান্নাঘর (rānnāghôr)}]{\bn{রান্নাঘর} (rānnāghôr) — a kitchen}

\label{lesson:BN-C508-rannaghor}



\begin{quote}
Before the new one: say the Bengali for a drawing room, then the Bengali for an attic.
\end{quote}

\subsection*{You'll want to know: \bn{রান্নাঘর}}
\textbf{\bn{রান্নাঘর}} — \emph{rānnāghôr} — \textquotedblleft{}a kitchen\textquotedblright{}.

\begin{quote}
\bn{দিদিমা} \bn{রান্নাঘরে} \bn{মাছ} \bn{ভাজছিলেন।}
\end{quote}

\begin{quote}
\emph{Didimā rānnāghôre māchh bhājchhilen.}
\end{quote}

\begin{quote}
My grandmother was frying fish in the kitchen.
\end{quote}

\begin{cousinweb}
\textbf{\bn{রান্নাঘর}} (\emph{rānnāghôr}) joins \bn{রান্না}, \textquotedblleft{}cooking\textquotedblright{} (CDIAL 10615), to \bn{ঘর}, \textquotedblleft{}a room\textquotedblright{} (CDIAL 4428): the cooking room, the everyday word for a kitchen.
\end{cousinweb}

\begin{grammarlens}[title={What you've built}]
One more word for describing what you lived and felt.
\end{grammarlens}

\subsection*{Your turn}
\begin{itemize}
\raggedright
  \item \emph{Say it:} \emph{rānnāghôr}
  \item \emph{Say it:} \emph{rānnāghôr}, once more
  \item \emph{Say it:} \emph{Didimā rānnāghôre māchh bhājchhilen.} — My grandmother was frying fish in the kitchen.
  \item \emph{Recall:} say the Bengali for a drawing room, then the Bengali for an attic, then say \emph{rānnāghôr} again
\end{itemize}

\begin{tcolorbox}[breakable,colback=teal!4,colframe=teal!35!black,title={Before you move on}]
What does \bn{রান্নাঘর} mean? (\textquotedblleft{}A kitchen\textquotedblright{}.) Say it once more.
\end{tcolorbox}
\section[\texorpdfstring{\bn{চৌকাঠ} (choukāṭh)}{চৌকাঠ (choukāṭh)}]{\bn{চৌকাঠ} (choukāṭh) — a door frame; a threshold}

\label{lesson:BN-C508-choukath}



\begin{quote}
Before the new one: say the Bengali for an attic, then the Bengali for a kitchen.
\end{quote}

\subsection*{You'll want to know: \bn{চৌকাঠ}}
\textbf{\bn{চৌকাঠ}} — \emph{choukāṭh} — \textquotedblleft{}a door frame; a threshold\textquotedblright{}.

\begin{quote}
\bn{মিস্ত্রি} \bn{দরজার} \bn{নতুন} \bn{চৌকাঠ} \bn{বানাল।}
\end{quote}

\begin{quote}
\emph{Mistri dôrjār nôtun choukāṭh bānālô.}
\end{quote}

\begin{quote}
The workman made a new frame for the door.
\end{quote}

Also the wooden sill at its foot: \bn{চৌকাঠে} \bn{বসা}, to sit on the threshold.

\begin{cousinweb}
\textbf{\bn{চৌকাঠ}} (\emph{choukāṭh}) comes from a reconstructed \emph{catuṣkāṣṭha}, \textquotedblleft{}a rectangular frame\textquotedblright{}, made of \emph{catur}, \textquotedblleft{}four\textquotedblright{}, and \emph{kāṣṭha}, \textquotedblleft{}wood\textquotedblright{}; CDIAL lists Bengali \emph{caukāṭh}, \textquotedblleft{}a wooden door frame\textquotedblright{} (CDIAL 4631).
\end{cousinweb}

\begin{grammarlens}[title={What you've built}]
One more word for describing what you lived and felt.
\end{grammarlens}

\subsection*{Your turn}
\begin{itemize}
\raggedright
  \item \emph{Say it:} \emph{choukāṭh}
  \item \emph{Say it:} \emph{choukāṭh}, once more
  \item \emph{Say it:} \emph{Mistri dôrjār nôtun choukāṭh bānālô.} — The workman made a new frame for the door.
  \item \emph{Recall:} say the Bengali for an attic, then the Bengali for a kitchen, then say \emph{choukāṭh} again
\end{itemize}

\begin{tcolorbox}[breakable,colback=teal!4,colframe=teal!35!black,title={Before you move on}]
What does \bn{চৌকাঠ} mean? (\textquotedblleft{}A door frame; a threshold\textquotedblright{}.) Say it once more.
\end{tcolorbox}
\section[\texorpdfstring{\bn{টাঙানো} (ṭāngāno)}{টাঙানো (ṭāngāno)}]{\bn{টাঙানো} (ṭāngāno) — to hang up (on a wall)}

\label{lesson:BN-C508-tangano}



\begin{quote}
Before the new one: say the Bengali for a kitchen, then the Bengali for a door frame.
\end{quote}

\subsection*{You'll want to know: \bn{টাঙানো}}
\textbf{\bn{টাঙানো}} — \emph{ṭāngāno} — \textquotedblleft{}to hang up (on a wall)\textquotedblright{}.

\begin{quote}
\bn{আমি} \bn{দেওয়ালে} \bn{দাদুর} \bn{একটা} \bn{ছবি} \bn{টাঙালাম।}
\end{quote}

\begin{quote}
\emph{Āmi deoyāle dādur ekṭā chhôbi ṭāngālām.}
\end{quote}

\begin{quote}
I hung a picture of my grandfather on the wall.
\end{quote}

Past (he/she …): \bn{তিনি} \bn{ছবিটা} \bn{টাঙালেন}, \emph{tini chhôbiṭā ṭāngālen}.

\begin{cousinweb}
\textbf{\bn{টাঙানো}} (\emph{ṭāngāno}) comes from a reconstructed root \emph{ṭaṅg}, \textquotedblleft{}hang\textquotedblright{}; CDIAL lists Bengali \emph{ṭāṅgā}, \textquotedblleft{}to hang up\textquotedblright{} (CDIAL 5436).
\end{cousinweb}

\begin{grammarlens}[title={What you've built}]
One more word for describing what you lived and felt.
\end{grammarlens}

\subsection*{Your turn}
\begin{itemize}
\raggedright
  \item \emph{Say it:} \emph{ṭāngāno}
  \item \emph{Say it:} \emph{ṭāngāno}, once more
  \item \emph{Say it:} \emph{Āmi deoyāle dādur ekṭā chhôbi ṭāngālām.} — I hung a picture of my grandfather on the wall.
  \item \emph{Recall:} say the Bengali for a kitchen, then the Bengali for a door frame, then say \emph{ṭāngāno} again
\end{itemize}

\begin{tcolorbox}[breakable,colback=teal!4,colframe=teal!35!black,title={Before you move on}]
What does \bn{টাঙানো} mean? (\textquotedblleft{}To hang up (on a wall)\textquotedblright{}.) Say it once more.
\end{tcolorbox}
